\documentclass[11pt]{amsart}

\usepackage{amsmath,amssymb,amsthm}
\usepackage{mathtools}
\usepackage[T1]{fontenc}
\usepackage{stmaryrd} % \llbracket \rrbracket
\usepackage[hidelinks]{hyperref}
\usepackage{enumitem}

% ---- theorem environments ----
\theoremstyle{plain}
\newtheorem{theorem}{Theorem}[section]
\newtheorem{proposition}[theorem]{Proposition}
\newtheorem{lemma}[theorem]{Lemma}
\newtheorem{corollary}[theorem]{Corollary}
\theoremstyle{definition}
\newtheorem{definition}[theorem]{Definition}
\newtheorem{example}[theorem]{Example}
\newtheorem{remark}[theorem]{Remark}
\newtheorem{question}[theorem]{Question}
\newtheorem{conjecture}[theorem]{Conjecture}
\newtheorem{problem}[theorem]{Problem}

% ---- notation ----
\newcommand{\Set}{\mathbf{Set}}
\newcommand{\Cont}{\mathbf{Cont}}
\newcommand{\Fam}{\mathrm{Fam}}
\newcommand{\Kl}{\mathrm{Kl}}
\newcommand{\Nat}{\mathrm{Nat}}
\newcommand{\Ran}{\mathrm{Ran}}
\newcommand{\op}{\mathrm{op}}
\newcommand{\ext}[1]{\llbracket #1 \rrbracket}
\newcommand{\extM}[1]{\llbracket #1 \rrbracket_{M}}
\newcommand{\id}{\mathrm{id}}
\newcommand{\Id}{\mathrm{Id}}
\newcommand{\ceil}[1]{\lceil #1 \rceil}
\newcommand{\tri}{\triangleleft}
\newcommand{\triM}{\triangleleft_{M}}
\newcommand{\y}{\mathsf{y}}
\newcommand{\D}{\mathsf{D}}
\newcommand{\Maybe}{\mathsf{Maybe}}
\newcommand{\PP}{\mathcal{P}}
\newcommand{\supp}{\mathrm{supp}}
\DeclareMathOperator{\im}{im}

\begin{document}

\title[Codensity and polynomiality of effectful containers]
{Two invariants of an effect monad: \\ codensity governs fullness, polynomiality governs composition}

\author{MacBeth}
\address{Kodamai}
\email{}
\date{6 September 2026}

\begin{abstract}
Fix a monad $M$ on $\Set$. An \emph{$M$-container} is an ordinary container $(S,P)$ whose
positions are read out in the Kleisli category of $M$; its extension is the endofunctor
$\extM{S,P}(X)=\sum_{s\in S}\Kl(M)(P_s,X)$, the effectful reading of a data type whose
observations may fail, branch, consult a context, or return a distribution. We prove three
things about the extension functor $\extM{-}\colon M\text{-}\mathbf{Cont}\to[\Set,\Set]$.
First, $M$-containers are exactly $\Fam(\Kl(M)^{\op})$ and the extension factors as
$\extM{-}=\ext{-}\circ M$: an $M$-container functor is an ordinary polynomial functor
post-composed with the effect monad. Second, $\extM{-}$ is \emph{always faithful}, and it is
\emph{full} if and only if $M$ is \emph{codense} --- the canonical map $M\Rightarrow\Ran_M M$
into its own codensity monad is an isomorphism --- and each functor $\Set(A,M{-})$ is connected;
consequently probability distributions keep fullness while error, partiality, read-only context
and nondeterminism destroy it. Third, the essential image of $\extM{-}$ is closed under
composition \emph{if} $M$ is a \emph{polynomial} functor --- sufficiency for all $M$, with
$p\triangleleft_M q=p\triangleleft\lceil M\rceil\triangleleft q$ --- and the converse holds precisely
within the \emph{fully symmetric} world. Composition-closure forces $M$ polynomial for every analytic
monad with $M\emptyset=\emptyset$ (by \emph{right-cancellation of plethysm} in the cycle-index ring, at
any degree or nesting depth, covering free magmas and free operad monads), for every analytic monad of
bounded \emph{wreath depth} (by an imprimitivity-depth invariant subsuming the multiset and
symmetric-power monads), and for every free commutative unital magma monad (by a correlated block-swap
no polynomial layer can produce); but it \emph{fails} for partially symmetric unital operads --- the free
algebra monad of a ternary operation with cyclic ($C_3$) slot-symmetry and an absorbed unit is
non-polynomial yet has composition-closed image ($\extM{p}\circ\extM{q}\cong\extM{p\triangleleft L}$
with $L$ the list monad). The dividing line is the commutativity of the binary unit-composite
$\mu(x,y)=\omega(x,y,e,\dots)$: commutative (fully symmetric) forces the block-swap and hence the
converse; non-commutative permits closure without polynomiality. Thus \emph{within the symmetric world
of the applied effect monads --- multiset, distribution, powerset --- composition $\iff$ flatness of the
species $\iff$ polynomiality}, with the $C_3$ operad marking the exact boundary. The affine-plus-commutative
folklore criterion is false ($\D$ is affine and commutative but does not compose; $\mathsf{Maybe}$ is
neither yet does). Fullness and composition are thus governed by two independent Kan-theoretic invariants
of $M$ that trade oppositely: probability is faithful but not composable, error is composable but not faithful.
\end{abstract}

\maketitle

\begin{center}\small
\fbox{\begin{minipage}{0.92\textwidth}\centering
\emph{Author:} MacBeth (Kodamai).\qquad\emph{Date:} 6 September 2026.\qquad
\emph{For review by:} Rick.\\[2pt]
\emph{Corresponds to} \texttt{scotmacbeth/work-in-progress} \emph{commit}
\texttt{PENDING}.\\[2pt]
\textbf{Draft --- work in progress; not yet neighbour-reviewed. This revision supersedes commit
\texttt{842db92}: Theorem~\ref{thm:3}'s necessity direction is now settled to be \emph{conditional}.
The last corner (formerly an open Problem, ``$U\circ U\cong\ext r\circ U$?'') is closed --- necessity
\emph{holds} at the free commutative unital magma $U$ (\S\ref{sec:unital}) --- but a partially symmetric
sibling, the $C_3$-ternary unital operad, \emph{refutes} necessity outright: its non-polynomial monad has
composition-closed image (Theorem~\ref{thm:C3}). The corrected theorem states sufficiency
unconditionally and necessity within the fully symmetric world, with commutativity of the binary
unit-composite as the exact dividing line (\S\ref{sec:unital}, Theorem~\ref{thm:3final}).}
\end{minipage}}
\end{center}

\section{Introduction}\label{sec:intro}

\subsection*{Containers, and what happens when the positions become effectful}
A \emph{container} $(S,P)$ --- a set $S$ of shapes together with a set $P_s$ of positions over
each shape --- presents a datatype through its extension, the endofunctor
\[
  \ext{S,P}(X)\;=\;\sum_{s\in S}\Set(P_s,X)\;=\;\sum_{s\in S}X^{P_s},
\]
which sends $X$ to ``a shape together with an $X$-labelling of its positions'' \cite{GambinoKock,NiuSpivak}.
The extensions of containers are exactly the \emph{polynomial functors}, and
$\ext{-}$ is fully faithful: a natural transformation between two such functors is precisely a
container morphism. This tight correspondence --- a syntax of shapes and positions on one side, a
class of functors on the other --- is what makes containers a usable calculus for datatypes.

Now suppose the positions are not read out purely but \emph{through an effect}. Reading a position
might fail, or return several answers, or consult an environment, or produce a probability
distribution over answers. All of these are captured by a monad $M$ on $\Set$: an effectful read
of a $P_s$-position landscape into $X$ is a Kleisli map $P_s\to X$ in $\Kl(M)$, i.e.\ a function
$P_s\to MX$. This yields the \emph{$M$-container} extension
\[
  \extM{S,P}(X)\;=\;\sum_{s\in S}\Kl(M)(P_s,X)\;=\;\sum_{s\in S}(MX)^{P_s}.
\]
The passage $\Cont\to M\text{-}\Cont$ that reinterprets a container's positions in $\Kl(M)$ is a
change of base of exactly the kind one wants for modelling effectful data structures and effectful
agents: an $M$-container is a device that, presented with a request, chooses a shape and then reads
its positions with the ambient effect. When $M=\Maybe$, an ``agent may decline'' a position; when
$M=\D$ (finite distributions), it answers probabilistically.

\subsection*{Two questions, and the temptation to conflate them}
For the pure calculus, the extension functor is fully faithful and strong monoidal: it names
functors faithfully, and those functors compose. Two properties, one embedding. When we move to
$M$-containers, both properties become genuine questions:
\begin{enumerate}[label=(Q\arabic*)]
\item\label{q:full} \textbf{Faithful naming.} Is $\extM{-}\colon M\text{-}\Cont\to[\Set,\Set]$
  full and faithful? That is, does every natural transformation between $M$-container functors come
  from a (unique) $M$-container morphism?
\item\label{q:comp} \textbf{Composition.} Is the essential image of $\extM{-}$ closed under functor
  composition? That is, if $F$ and $G$ are $M$-container functors, is $F\circ G$ one too, presented
  by some container built from theirs?
\end{enumerate}
It is tempting to expect that these rise and fall together --- that ``$M$-containers behave well''
is a single condition on $M$. The main discovery of this note is that \ref{q:full} and \ref{q:comp}
are controlled by \emph{two different} invariants of $M$, and that across the standard effects they
trade against each other.

\subsection*{Results}
Throughout, $M=(M,\eta,\mu)$ is a monad on $\Set$, $\Kl(M)$ its Kleisli category, and
$\ext{-}\colon\Cont\to[\Set,\Set]$ the ordinary (polynomial) extension.

The first theorem identifies the players and provides the lever for everything else.

\begin{theorem}[Identification and factorisation; \S\ref{sec:identify}]\label{thm:1}
The category of $M$-containers is $\Fam(\Kl(M)^{\op})$: an object is an ordinary container
$(S,P\colon S\to\Set)$, and a morphism $(S,P)\to(S',P')$ is a shape map $f\colon S\to S'$ together
with, for each $s$, a Kleisli map $\rho_s\in\Kl(M)(P'_{fs},P_s)=\Set(P'_{fs},MP_s)$. Moreover, as
functors $\Set\to\Set$ and naturally in $(S,P)$,
\[
  \extM{S,P}\;=\;\ext{S,P}\circ M .
\]
\end{theorem}

So an $M$-container functor is an \emph{ordinary polynomial functor pre-composed with the effect
monad $M$}. Neil Ghani's ``change of base'' $\Cont\to M\text{-}\Cont$ is precomposition with the
Kleisli inclusion $F\colon\Set\to\Kl(M)$, and every structural question about $\extM{-}$ becomes a
question about the single trailing $M$.

The second theorem answers \ref{q:full}. Recall that the \emph{codensity monad} of a functor
$M\colon\Set\to\Set$ is the right Kan extension $\Ran_M M$; there is a canonical monad map
$\iota_M\colon M\Rightarrow\Ran_M M$, and $M$ is called \emph{codense} when $\iota_M$ is an
isomorphism (Definition~\ref{def:codense}). Call $M$ \emph{affine} if $M1\cong 1$.

\begin{theorem}[Faithfulness always; fullness $=$ codensity; \S\ref{sec:full}]\label{thm:2}
The extension $\extM{-}\colon M\text{-}\Cont\to[\Set,\Set]$ is \emph{always faithful}. It is
\emph{full} --- hence fully faithful --- if and only if
\begin{enumerate}[label=\textup{(\roman*)}]
\item each functor $\Set(A,M{-})$ is connected, and
\item $M$ is codense.
\end{enumerate}
If $M$ is affine then \textup{(i)} holds; hence \emph{$M$ affine and codense $\Rightarrow$ fully
faithful}. The finite distribution monad $\D$ is affine and codense, so $\extM[\D]{-}$ is fully
faithful \textup{(}modulo the rigidity Lemma~\ref{lem:Drigid}\textup{)}; whereas $\Maybe$,
exceptions, the writer $E\times(-)$, the reader $(-)^R$ and nonempty powerset all fail codensity and
so lose fullness.
\end{theorem}

The correct invariant is the codensity monad, \emph{not} preservation of coproducts. The natural
first guess --- ``$\extM{-}$ is fully faithful iff $M$ preserves coproducts'' --- is wrong for the
target $[\Set,\Set]$: it is the criterion for a coarser, Kleisli-enriched target, and it gives the
wrong verdict on $\D$. The writer monad $E\times(-)$ preserves coproducts yet fails fullness for
$|E|\ge 2$, a clean refutation (\S\ref{sec:correction}).

The third theorem answers \ref{q:comp}. Call $M$ \emph{polynomial} if $M\cong\sum_{i\in I}(-)^{B_i}$
for some family $(B_i)_{i\in I}$; equivalently $M=\ext{\ceil M}$ for a container $\ceil M$.

\begin{theorem}[Sufficiency of polynomiality; conditional converse; \S\ref{sec:comp}]\label{thm:3}
\emph{(Sufficiency, unconditional.)} If $M$ is polynomial, then for all $M$-containers $p,q$ the
ordinary container
\[
  p\triangleleft_M q\;:=\;p\triangleleft\ceil M\triangleleft q
\]
satisfies $\extM{p\triangleleft_M q}\cong\extM{p}\circ\extM{q}$, naturally in $X$. The operation
$\triangleleft_M$ is associative and has a unit up to ordinary container isomorphism only when
$M=\Id$; so the essential image of $\extM{-}$ is closed under composition.
\emph{(Necessity, conditional.)} The converse --- composition-closure $\Rightarrow$ $M$ polynomial ---
holds for the following classes, but not in general. It holds if $M$ grows super-polynomially (a
cardinality law excludes powerset, nonempty powerset, distributions, ultrafilter, continuation); if $M$
is analytic with $M\emptyset=\emptyset$ (right-cancellation of plethysm in the cycle-index ring,
Theorem~\ref{thm:P}, at any degree or nesting depth, covering free magmas and free operad monads); if
$M$ is analytic of bounded \emph{wreath depth} (an imprimitivity-depth invariant, Theorem~\ref{thm:Sprime},
subsuming the multiset and symmetric-power monads); and if $M$ is a free commutative unital magma monad
(Theorem~\ref{thm:family}, including the corner case $U$ formerly left open). \emph{But the converse
fails} for partially symmetric unital operads: the free algebra monad of a single ternary operation with
cyclic ($C_3$) slot-symmetry and an absorbed unit is analytic and non-polynomial, yet its essential image
\emph{is} closed under composition (Theorem~\ref{thm:C3}). The exact dividing line is the commutativity of
the binary unit-composite $\mu(x,y)=\omega(x,y,e,\dots)$: when $\mu$ is commutative (full slot-symmetry),
$M\circ M$ realises a correlated block-swap that no polynomial layer $\ext r\circ M$ can manufacture, so
necessity holds; when $\mu$ is non-commutative, an ``only three identical'' rigidity lets the image close
without polynomiality. Consequently \textup{(}Theorem~\ref{thm:3final}\textup{)} \emph{within the fully
symmetric world --- which contains all the applied effect monads: multiset, distribution, powerset ---
composition-closure $\iff$ flatness of the species $\iff$ polynomiality}, and the $C_3$ operad marks the
precise boundary at which the symmetry hypothesis becomes load-bearing.
\end{theorem}

Combining the two theorems gives the punchline.

\begin{corollary}[Independence and trade-off]\label{cor:tradeoff}
Fullness \textup{(codensity)} and composability \textup{(polynomiality)} are independent invariants
of $M$. The writer monad $E\times(-)$ is polynomial, hence composes, but is not codense, hence not
full. The distribution monad $\D$ is affine and codense, hence full, but is not polynomial, hence
does not compose. \emph{Probability keeps faithful naming but loses composition; error keeps
composition but loses faithful naming.}
\end{corollary}

\subsection*{Method}
Everything runs off the factorisation $\extM{-}=\ext{-}\circ M$ of Theorem~\ref{thm:1}. For
faithfulness, the Kleisli unit $\eta$ is a section, so container data can be read back off the
induced natural transformation. For fullness, precomposition $(-)\circ M$ has a right adjoint
$\Ran_M$, and every hom-set $\Nat(\Set(A,M{-}),H)$ is a value of a right Kan extension along $M$;
the fullness obstruction is exactly the failure of the unit $M\Rightarrow\Ran_M M$ to be invertible.
For composition, two applications of the factorisation put every composite $\extM{p}\circ\extM{q}$
in the form $G\circ M$ with $G=\ext{p}\circ M\circ\ext{q}$; this is again an $M$-extension precisely
when $G$ is polynomial, and then AAG strong monoidality delivers the splice $p\triangleleft\ceil
M\triangleleft q$. Sufficiency is immediate. The converse is where the structure lives, and it is
genuinely partial. The trailing $M$ cannot be cancelled at the level of connected-limit preservation
--- Proposition~\ref{prop:inert} shows the single self-composite $M\circ M\cong\ext r\circ M$ carries
no information about $M$ beyond the retract reduction Lemma~\ref{lem:R0} --- so the converse needs a
\emph{finer, faithful} invariant. The master lemma is that a \emph{polynomial} outer layer is
rigid on positions: every point-stabilizer of $(\ext r\circ M)[m]$ is a \emph{block-fixing product}
of single-$M$-element stabilizers (Proposition~\ref{prop:meta}, ``a function is fixed iff fixed
pointwise''). Three invariants then measure how far $M\circ M$ escapes this constraint: a cardinality
count for super-polynomial growth; an imprimitivity-\emph{wreath-depth} statistic $h$
(Definition~\ref{def:depth}) for bounded-depth analytic monads; and right-cancellation of plethysm on
the cycle-index series (Lemma~\ref{lem:cancel}) for analytic $M$ with $M\emptyset=\emptyset$, reaching
the unbounded-depth monads no stabilizer count can touch. The residual after these --- analytic,
$M\emptyset$ nonempty-finite, unbounded wreath depth --- is where symmetry becomes decisive. If the
monad's operations are \emph{fully} symmetric with an absorbed unit, then $M\circ M$ builds a
\emph{rigid} composite block from the unit and \emph{swaps} two equal copies, a correlated symmetry
$\Delta_{C_2}\notin\mathcal F$ that no block-fixing product contains, and necessity holds
(Theorem~\ref{thm:family}). If instead a binary composite is non-commutative --- as for the $C_3$
operad --- an ``only three identical'' rigidity (Lemma~\ref{lem:c3}) shows every stabilizer is already
a single-tree automorphism group, the image closes, and necessity fails (Theorem~\ref{thm:C3}).

\subsection*{Why it matters}
Beyond the pleasing tension of Corollary~\ref{cor:tradeoff}, the result sharpens the compositional
picture of effectful agents. Modelling ``an agent may decline'' as a $\Maybe$-container, one keeps
composition (agents chain into agents) at the price of faithfulness (distinct declining strategies
can induce the same functor). Modelling ``an agent answers probabilistically'' as a
$\D$-container, one keeps faithful naming (strategies are recoverable from behaviour) but loses
strict composition (a chain of probabilistic agents is not itself presented by a single
probabilistic container --- the marginal-to-joint gap). The invariant that decides each is now
explicit and computable: codensity for the first question, polynomiality for the second.

\subsection*{Provenance and honesty}
Theorems~\ref{thm:1} and \ref{thm:2} (faithfulness, the codensity criterion, the affine lemma, the
examples) are proved. Theorem~\ref{thm:3} is proved in its sufficiency direction unconditionally, and in
its necessity direction for each of the four classes named: the cardinality law disposes of the
super-polynomial monads; plethysm right-cancellation (Lemma~\ref{lem:cancel}, Theorem~\ref{thm:P})
disposes of every analytic monad with $M\emptyset=\emptyset$, at unbounded degree and depth; the
wreath-depth invariant (Definition~\ref{def:depth}, Theorem~\ref{thm:Sprime}) disposes of every
bounded-wreath-depth analytic monad, subsuming the multiset and symmetric-power monads; and the
correlated-swap witness (Theorem~\ref{thm:family}) disposes of every free commutative unital magma monad,
including the corner case $U$. The failure of necessity for the $C_3$ operad (Theorem~\ref{thm:C3}) is
likewise proved: the isomorphism $M\circ M\cong L\circ M$ is established by matching molecule multisets
(the ``only three identical'' rigidity, Lemma~\ref{lem:c3}); the small cases $m\le4$ are confirmatory
computation, not the proof. The cancellation and block-fixing engines are elementary and computationally
verified. Theorems~\ref{thm:P}, \ref{thm:Sprime}, \ref{thm:family} and \ref{thm:C3} invoke two standard
facts of combinatorial-species theory --- Joyal's full faithfulness of the species-to-analytic-functor
embedding, and the identity of species substitution with plethysm of cycle-index series (Joyal;
Bergeron--Labelle--Leroux; Gambino--Kock for the polynomial-$\iff$-flat criterion) --- which we cite as
background rather than re-derive. What remains genuinely open is the necessity direction in the residual
outside the free commutative unital magmas: analytic $M$ with $M\emptyset$ nonempty-finite, unbounded
wreath depth, and \emph{not} fully symmetric (Problem~\ref{prob:residual}). Full faithfulness for $\D$
rests on a standard rigidity fact, $\Nat(\D,\D)=\{\id\}$, which we isolate as Lemma~\ref{lem:Drigid} and
assume rather than reprove. These boundaries are marked at each occurrence.

\subsection*{Outline}
Section~\ref{sec:prelim} fixes notation: monads and Kleisli categories, the container extension and
its composition monoidal structure, families over an opposite category, and right Kan extensions and
codensity monads. Section~\ref{sec:identify} proves Theorem~\ref{thm:1} and the always-faithfulness.
Section~\ref{sec:full} proves the codensity criterion (Theorem~\ref{thm:2}), works the examples, and
explains why coproduct preservation is the wrong criterion. Section~\ref{sec:comp} proves
Theorem~\ref{thm:3}: sufficiency and non-unitality (\S\ref{sec:suff}); the cardinality law
(\S\ref{sec:card}); the block-fixing meta-theorem and stabilizer invariant (\S\ref{sec:stab}); the
wreath-depth invariant and Theorem~S$'$ (\S\ref{sec:depth}); the plethysm right-cancellation closing the
analytic converse for $M\emptyset=\emptyset$ (\S\ref{sec:cancel}); and the unital corner
(\S\ref{sec:unital}), where full symmetry gives necessity via a correlated block-swap while the $C_3$
operad refutes it, pinning commutativity of the binary unit-composite as the exact dividing line.
Section~\ref{sec:conc} draws the trade-off and lists open problems.

\section{Preliminaries}\label{sec:prelim}

We work over $\Set$. All the constructions below have counterparts over a general base along the
lines of $\Fam(\mathcal C^{\op})$ \cite{DJN}, but the codensity and polynomiality phenomena are
already visible, and cleanest, over $\Set$.

\subsection{Monads and Kleisli categories}
A monad $M=(M,\eta,\mu)$ on $\Set$ has unit $\eta\colon\Id\Rightarrow M$ and multiplication
$\mu\colon MM\Rightarrow M$. Its \emph{Kleisli category} $\Kl(M)$ has sets as objects,
$\Kl(M)(A,B)=\Set(A,MB)$, identity $\eta_A$, and composition $g\bullet f=\mu\circ Mg\circ f$. The
free functor $F\colon\Set\to\Kl(M)$ is the identity on objects and sends $g\colon A\to B$ to
$\eta_B\circ g$; it is faithful and bijective on objects, is a left adjoint (so preserves colimits),
and is \emph{not} full for nontrivial $M$. We call $M$ \emph{affine} if $M1\cong1$, and
\emph{commutative} if its two canonical strengths $MA\times MB\to M(A\times B)$ agree.

\subsection{Containers and the polynomial extension}
A \emph{container} is a pair $(S,P)$ with $S$ a set and $P\colon S\to\Set$ (write $P_s$ for the
fibre). The category $\Cont$ has as morphisms $(S,P)\to(S',P')$ the pairs $(f,\varphi)$ of a shape
map $f\colon S\to S'$ and a reindexing $\varphi_s\colon P'_{fs}\to P_s$; equivalently
$\Cont=\Fam(\Set^{\op})$ (see below). The \emph{extension} is
\[
  \ext{S,P}(X)=\sum_{s\in S}\Set(P_s,X),\qquad \ext{S,P}(g)=\textstyle\sum_s(g\circ-).
\]
The functor $\ext{-}\colon\Cont\to[\Set,\Set]$ is fully faithful, with essential image the
\emph{polynomial functors} $\sum_{i\in I}(-)^{B_i}$ \cite[\S1]{GambinoKock},\cite{NiuSpivak}. We say
$M$ is \emph{polynomial} when $M$ lies in this image, and write $\ceil M$ for a container with
$\ext{\ceil M}=M$.

$\Cont$ carries the \emph{composition} monoidal structure $(\Cont,\triangleleft,\y)$, where $\y=(1,
1\mapsto1)$ is the identity container ($\ext{\y}=\Id$) and $\triangleleft$ is defined so that
$\ext{-}$ is strong monoidal:
\begin{equation}\label{eq:AAGmon}
  \ext{c\triangleleft d}\;\cong\;\ext{c}\circ\ext{d},\qquad \ext{\y}=\Id,
\end{equation}
naturally, and $\triangleleft$ is associative up to coherent isomorphism
\cite[Thm.~1.12]{GambinoKock},\cite{NiuSpivak}. Directed containers --- containers whose
$\triangleleft$-structure is a comonoid --- are the categories internal to this calculus
\cite{AhmanUustaluDCC}.

\subsection{Families over an opposite category}
For a category $\mathcal C$, the category $\Fam(\mathcal C)$ has objects $(S,X\colon S\to\mathrm{Ob}\,
\mathcal C)$ and morphisms $(S,X)\to(S',X')$ given by a function $f\colon S\to S'$ and a family of
$\mathcal C$-maps $X_s\to X'_{fs}$. Thus $\Fam(\mathcal C^{\op})$ has the same objects but morphisms
whose component maps run $X'_{fs}\to X_s$ in $\mathcal C$. In particular $\Cont=\Fam(\Set^{\op})$.

\subsection{Right Kan extensions and the codensity monad}
Precomposition $M^{*}=(-)\circ M\colon[\Set,\Set]\to[\Set,\Set]$ has a right adjoint, the right Kan
extension $\Ran_M\colon[\Set,\Set]\to[\Set,\Set]$ along $M$. The defining adjunction gives, for
$H\colon\Set\to\Set$ and any representable $\Set(A,-)$,
\begin{equation}\label{eq:kan}
\begin{aligned}
  \Nat\bigl(\Set(A,M{-}),\,H\bigr)
  &=\Nat\bigl(\Set(A,-)\circ M,\,H\bigr)\\
  &=\Nat\bigl(\Set(A,-),\,\Ran_M H\bigr)=(\Ran_M H)(A),
\end{aligned}
\end{equation}
using the Yoneda lemma in the last step. Applying this with $H=M$ produces the \emph{codensity
monad} of $M$.

\begin{definition}\label{def:codense}
The \emph{codensity monad} of $M$ is $\Phi:=\Ran_M M$; by \eqref{eq:kan},
$\Phi(A)=\Nat(\Set(A,M{-}),M)$. The unit $\eta$ of the adjunction $M^{*}\dashv\Ran_M$, evaluated at
$M$, is a natural monad map $\iota_M\colon M\Rightarrow\Phi$, explicitly
$\iota_{M,A}(m)=[\,k\mapsto\mu_A(Mk\,(m))\,]$. We call $M$ \emph{codense} when $\iota_M$ is an
isomorphism.
\end{definition}

\begin{definition}\label{def:connected}
A functor $G\colon\Set\to\Set$ is \emph{connected} if $\Nat(G,-)$ preserves coproducts:
$\Nat(G,\sum_j H_j)=\sum_j\Nat(G,H_j)$ for every family $(H_j)$. Equivalently, every natural
transformation out of $G$ into a coproduct factors through a single summand.
\end{definition}

\section{The identification, and always-faithfulness}\label{sec:identify}

\begin{definition}\label{def:mcont}
An \emph{$M$-container} is a container $(S,P)$; its \emph{$M$-extension} is
\[
  \extM{S,P}(X)=\sum_{s\in S}\Kl(M)(P_s,X)=\sum_{s\in S}\Set(P_s,MX),
\]
with action $\extM{S,P}(g)=\sum_s(Mg\circ-)$ on a map $g\colon X\to X'$.
A morphism $(S,P)\to(S',P')$ of $M$-containers is a shape map $f\colon S\to S'$ together with, for
each $s$, a Kleisli map $\rho_s\in\Kl(M)(P'_{fs},P_s)=\Set(P'_{fs},MP_s)$; composition uses
$\bullet$ in $\Kl(M)$.
\end{definition}

\begin{proof}[Proof of Theorem~\ref{thm:1}]
\emph{(a) $M\text{-}\Cont=\Fam(\Kl(M)^{\op})$.} Objects of $\Fam(\Kl(M)^{\op})$ are pairs
$(S,P\colon S\to\mathrm{Ob}\,\Kl(M))$, and $\mathrm{Ob}\,\Kl(M)=\Set$, so they are exactly
containers. A morphism $(S,P)\to(S',P')$ is $f\colon S\to S'$ together with, per $s$, a
$\Kl(M)^{\op}$-map $P_s\to P'_{fs}$, i.e.\ a $\Kl(M)$-map $P'_{fs}\to P_s$, i.e.\ a function
$P'_{fs}\to MP_s$ --- exactly Definition~\ref{def:mcont}. Identities match ($\Fam$'s identity uses
$\id$ in $\Kl(M)^{\op}$, i.e.\ $\eta$, the $M$-container identity), and composition matches
($\Fam$ composes shapes in $\Set$ and positions by $\bullet$ in the opposite order). The two
categories are equal on objects, hom-sets, identities and composition.

\emph{(b) $\extM{S,P}=\ext{S,P}\circ M$.} On objects,
$\ext{S,P}(MX)=\sum_s\Set(P_s,MX)=\extM{S,P}(X)$; on $g\colon X\to X'$ both sides act by
$\sum_s(-\circ Mg)$. For naturality in $(S,P)$, a morphism $(f,\rho)$ induces
$\alpha_X(s,k)=(fs,\,k\bullet\rho_s)=(fs,\,\mu_X\circ Mk\circ\rho_s)$, which is precisely $\ext{-}\circ
M$ applied to the underlying shape/position data.
\end{proof}

Part (b) is the lever: \emph{an $M$-container functor is a polynomial functor with a trailing $M$.}

\begin{proposition}[Always faithful]\label{prop:faithful}
$\extM{-}\colon M\text{-}\Cont\to[\Set,\Set]$ is faithful.
\end{proposition}

\begin{proof}
Suppose $(f,\rho)$ and $(\tilde f,\tilde\rho)$ induce the same natural transformation $\alpha$.
Evaluate at $X=P_s$ on the element $(s,\eta_{P_s})\in\extM{S,P}(P_s)$:
$\alpha_{P_s}(s,\eta_{P_s})=(fs,\,\eta_{P_s}\bullet\rho_s)=(fs,\rho_s)$, since $\eta$ is the Kleisli
identity. Equality of $\alpha$ therefore forces $fs=\tilde fs$ and $\rho_s=\tilde\rho_s$ for every
$s$. Hence $(f,\rho)=(\tilde f,\tilde\rho)$.
\end{proof}

\begin{remark}
Faithfulness uses only that $\eta$ is a section of the Kleisli identity; it holds for every monad.
This already separates $M$-containers from linear (vector-space) containers, whose extension is not
faithful because the analogue of $\eta$ is not a section.
\end{remark}

\section{Fullness is codensity}\label{sec:full}

Fix $M$-containers $(S,P)$ and $(S',P')$. By Theorem~\ref{thm:1}(b), $\extM{S,P}=p\circ M$ and
$\extM{S',P'}=p'\circ M$ with $p=\ext{S,P}$, $p'=\ext{S',P'}$. Since $p=\sum_s\Set(P_s,-)$ is a
coproduct of representables,
\begin{equation}\label{eq:homnat}
  \Nat(p\circ M,\;p'\circ M)=\prod_{s}\Nat\Bigl(\Set(P_s,M{-}),\;\sum_{s'}\Set(P'_{s'},M{-})\Bigr),
\end{equation}
while the $M$-container hom-set is $\prod_s\sum_{s'}\Set(P'_{s'},MP_s)$. Comparing, fullness
holds for all pairs if and only if, for every set $A$ (playing the role of $P_s$) and every family
$(P'_{s'})$, the map
\begin{equation}\label{eq:star}\tag{$\star$}
\begin{aligned}
  \sum_{s'}\Set(P'_{s'},MA)&\;\longrightarrow\;\Nat\Bigl(\Set(A,M{-}),\;\sum_{s'}\Set(P'_{s'},M{-})\Bigr),\\
  (s',\rho)&\;\mapsto\;[\,k\mapsto(s',k\bullet\rho)\,],
\end{aligned}
\end{equation}
is a bijection.

\subsection{The Kan engine and the split}
By \eqref{eq:kan} with $A$, every natural transformation out of $\Set(A,M{-})$ is a value of a right
Kan extension along $M$. Two features of $\Set$ let us split \eqref{eq:star} cleanly. First, a
single summand is a product of copies of $M$: $\Set(B,M{-})=\prod_{b\in B}M$, so
$\Nat(\Set(A,M{-}),\Set(B,M{-}))=\Nat(\Set(A,M{-}),M)^{B}=\Phi(A)^B$ by
Definition~\ref{def:codense}. Second, the coproduct over $s'$ is handled exactly when
$\Set(A,M{-})$ is connected. This yields:

\begin{proof}[Proof of Theorem~\ref{thm:2}]
Faithfulness is Proposition~\ref{prop:faithful}. For fullness we show \eqref{eq:star} is a bijection
for all $A,(P'_{s'})$ iff (i) and (ii) hold.

\emph{(ii) handles a single summand.} With one summand $B=P'_{s'}$, \eqref{eq:star} becomes
$\Set(B,MA)\to\Nat(\Set(A,M{-}),\Set(B,M{-}))=\Phi(A)^B$, which is $\bigl(MA\to\Phi(A)\bigr)^B$; and
the map is $\iota_{M,A}$ in each coordinate. So the single-summand case for all $B$ is a bijection
iff $\iota_{M,A}\colon MA\to\Phi(A)$ is a bijection for all $A$ --- that is, iff $M$ is codense.

\emph{(i) handles the coproduct.} Given codensity, the general \eqref{eq:star} is a bijection iff the
canonical comparison, with $G_{s'}=\Set(P'_{s'},M{-})$,
\[
  \textstyle\sum_{s'}\Nat(\Set(A,M{-}),G_{s'})\;\longrightarrow\;\Nat(\Set(A,M{-}),\sum_{s'}G_{s'})
\]
is a bijection --- that is, iff $\Set(A,M{-})$ is connected.

Hence \eqref{eq:star} for all data $\iff$ (i)$\wedge$(ii), proving the criterion.
\end{proof}

\begin{lemma}[Affine $\Rightarrow$ connected]\label{lem:affine}
If $M$ is affine then each $\Set(A,M{-})$ is connected; hence $M$ affine and codense $\Rightarrow$
$\extM{-}$ is fully faithful.
\end{lemma}

\begin{proof}
Since $M1\cong1$, $\Set(A,M1)=(M1)^A\cong1$: the functor $\Set(A,M{-})$ has a single global point.
Let $\alpha\colon\Set(A,M{-})\Rightarrow\sum_{s'}G_{s'}$. The unique point of $\Set(A,M1)$ maps to a
single summand $s'_0$ of $\sum_{s'}G_{s'}(1)$. For any $X$, the terminal map $!\colon X\to1$ gives a
commuting square; since $\sum_{s'}G_{s'}(!)$ preserves summands and the image of the global point
lies in $s'_0$, naturality forces $\alpha_X$ to land in the $s'_0$-summand for every $X$. Thus
$\alpha$ factors through $s'_0$, which is connectedness.
\end{proof}

\subsection{The examples}
For a polynomial monad $M=\sum_{i\in I}(-)^{B_i}$ the whole comparison is computable by
\eqref{eq:AAGmon}: $\Set(A,M{-})$ is again polynomial, with shapes $\varphi\colon A\to I$ and
positions $\sum_{a}B_{\varphi a}$, so that
\[
  |MA|=\sum_{i}|A|^{|B_i|},\qquad
  |\Phi(A)|=\prod_{\varphi\colon A\to I}\;\sum_{i\in I}\Bigl(\textstyle\sum_a
  |B_{\varphi a}|\Bigr)^{|B_i|}.
\]
Codensity is the equality $|MA|=|\Phi(A)|$ for all $A$. Exact computation gives:

\begin{center}
\renewcommand{\arraystretch}{1.2}
\begin{tabular}{lllcc}
\hline
$M$ & $M1$ & $|MA|$ vs $|\Phi(A)|$ ($A=0,1,2,3$) & codense? & full?\\
\hline
$\Id$ & $1$ & $0,1,2,3$ \; \textbf{=} \; $0,1,2,3$ & yes & \textbf{yes}\\
$\Maybe$ $(X{+}1)$ & $2$ & $1,2,3,4$ vs $1,2,\mathbf{12},\mathbf{864}$ & no & no\\
exception $X{+}2$ & $3$ & $2,3,4,5$ vs $2,\mathbf{12},\mathbf{5184},\dots$ & no & no\\
writer $2{\times}X$ & $2$ & $0,2,4,6$ vs $0,\mathbf{4},\mathbf{256},\dots$ & no & no\\
reader $X^{2}$ & $1$ & $0,1,4,9$ vs $0,\mathbf{4},\mathbf{16},\mathbf{36}$ & no & no\\
\hline
\end{tabular}
\end{center}

Beyond a low threshold $|\Phi(A)|>|MA|$ strictly in every nontrivial case: the surplus natural
transformations are ``case-analysis'' operations --- inspecting $\bot$, or duplicating a reader
argument --- that are natural for pure maps but not Kleisli-natural.

\begin{example}[Reader: affine but not codense]\label{ex:reader}
The reader monad $(-)^R$ is affine ($M1=1$, so (i) holds by Lemma~\ref{lem:affine}) yet fails (ii):
$\Phi(A)=(RA)^{R}\ne A^{R}=MA$ for $|R|\ge2$. Affineness alone does not give fullness; codensity is a
genuinely separate necessary condition.
\end{example}

\begin{example}[Distributions: affine and codense]\label{ex:D}
Let $\D$ be the finitely-supported distribution monad. It is affine ($\D1\cong1$). We claim it is
codense: $\Phi(A)=\Nat(\D^{A},\D)=\D A$, i.e.\ the natural $A$-ary operations on $\D$ are exactly the
affine (convex) combinations. Given $\alpha\colon\D^{A}\Rightarrow\D$ and a tuple
$(\mu_a)_a\in(\D X)^A$, put $K=\bigsqcup_a\supp(\mu_a)$, let $\nu_a\in\D K$ be the block-$a$ copy of
$\mu_a$, and $p\colon K\to X$ the projection; then $(\mu_a)=\D^{A}(p)((\nu_a))$, so
$\alpha_X((\mu_a))=\D p\,(\alpha_K((\nu_a)))$. Collapsing the blocks by $c\colon K\to A$ gives
$\D c\,(\alpha_K((\nu_a)))=\alpha_A((\delta_a))=:\omega\in\D A$. Reducing one block at a time (collapse
all but one block to a point) reduces to classifying $\Nat(\D,\D(-\sqcup F))$ for finite $F$: the
$F$-marginal is constant (it factors through $\D1=1$) and the remaining marginal is a natural
sub-probability-valued endofunctor of $\D$, hence a scalar $\lambda\cdot\sigma$. Thus
$\alpha_K((\nu_a))=\sum_a\omega_a\nu_a$ and $\alpha_X((\mu_a))=\sum_a\omega_a\mu_a$. So $\D$ is
codense, and by Lemma~\ref{lem:affine} \emph{$\extM[\D]{-}$ is fully faithful} --- modulo the base
rigidity Lemma~\ref{lem:Drigid}.
\end{example}

\begin{lemma}[Rigidity of $\D$; assumed]\label{lem:Drigid}
The only natural endomorphism of $\D$ is the identity, and every natural sub-probability-valued
endofunctor of $\D$ is a scalar multiple of the identity: $\Nat(\D,\D)=\{\id\}$.
\end{lemma}

This is a standard rigidity fact for the affine functor $\D$; we assume it rather than reprove it,
and we mark the full faithfulness of $\extM[\D]{-}$ as conditional on it. The reduction in
Example~\ref{ex:D} is unconditional.

\subsection{Why coproduct preservation is the wrong criterion}\label{sec:correction}
A natural first guess is that $\extM{-}$ is fully faithful iff $M$ preserves coproducts. This is
false for the target $[\Set,\Set]$, for a structural reason. Factor
\[
  \extM{-}\;=\;R\circ\ext{-}^{\Kl},\qquad R=(-\circ F),
\]
where $\ext{-}^{\Kl}\colon\Fam(\Kl(M)^{\op})\to[\Kl(M),\Kl(M)]$ is the free-coproduct-completion
extension (always fully faithful) and $R$ restricts along the free functor $F\colon\Set\to\Kl(M)$.
``$M$ preserves coproducts'' is the fullness criterion for the coarser target
$[\Kl(M),\Kl(M)]$ --- the connected-unit condition for $\ext{-}^{\Kl}$ --- and $\ext{-}^{\Kl}$ is
already fully faithful. All the genuine content is in $R$, i.e.\ in the gap between pure-naturality
and Kleisli-naturality, and that gap is measured by the codensity monad, not by coproduct
preservation.

The verdicts differ concretely. Coproduct preservation fails for $\D$, so the wrong criterion
predicts $\D$ fails fullness; the correct criterion (affine and codense) says $\D$ \emph{keeps}
fullness. Conversely the writer monad $E\times(-)$ \emph{preserves} coproducts, so the wrong
criterion predicts fullness, yet $\Phi(1)=|E|^{|E|}\ne|E|$ for $|E|\ge2$, so $E\times(-)$ is not
codense and $\extM{-}$ is not full. The writer monad is thus a clean refutation of the coproduct
criterion.

\begin{remark}[Corrected slogan]
Adding probabilistic positions preserves the on-the-nose faithful naming of the container extension;
adding error, partiality, read-only context or nondeterminism destroys it. The dividing invariant is
codensity of the effect monad, computed by the right Kan extension $\Ran_M M$.
\end{remark}

\section{Composition is polynomiality}\label{sec:comp}

We now turn to \ref{q:comp}. Throughout, $\ext{-}\colon\Cont\to[\Set,\Set]$ is the strong monoidal
extension of \eqref{eq:AAGmon}.

\subsection{The structural identity}
\begin{lemma}\label{lem:A0}
For all $M$-containers $p,q$, as functors $\Set\to\Set$ and naturally in $X$,
\[
  \extM{p}\circ\extM{q}\;=\;G\circ M,\qquad G:=\ext{p}\circ M\circ\ext{q}.
\]
\end{lemma}

\begin{proof}
Two applications of Theorem~\ref{thm:1}(b) and associativity of $\circ$:
\[
  (\extM{p}\circ\extM{q})(X)=\extM{p}\bigl(\ext{q}(MX)\bigr)=\ext{p}\bigl(M(\ext{q}(MX))\bigr)
  =(\ext{p}\circ M\circ\ext{q}\circ M)(X).\qedhere
\]
\end{proof}

So every composite of two $M$-extensions has a trailing $M$; it is again an $M$-extension exactly
when the inner $G$ is polynomial.

\subsection{Sufficiency}\label{sec:suff}
\begin{proposition}[Sufficiency]\label{prop:suff}
If $M$ is polynomial with container $\ceil M$, then $p\triangleleft_M q:=p\triangleleft\ceil
M\triangleleft q$ satisfies $\extM{p\triangleleft_M q}\cong\extM{p}\circ\extM{q}$ naturally in $X$.
The operation $\triangleleft_M$ on $\Cont$ is associative up to the $\triangleleft$-associator.
\end{proposition}

\begin{proof}
With $M=\ext{\ceil M}$, \eqref{eq:AAGmon} applied twice gives
\[
  G=\ext{p}\circ M\circ\ext{q}=\ext{p}\circ\ext{\ceil M}\circ\ext{q}\cong\ext{p\triangleleft\ceil
  M\triangleleft q}=\ext{r},\qquad r:=p\triangleleft\ceil M\triangleleft q,
\]
naturally, the two bracketings agreeing by associativity of $\triangleleft$. Whiskering by $M$ on
the right and invoking Lemma~\ref{lem:A0},
\[
  \extM{p}\circ\extM{q}=G\circ M\cong\ext{r}\circ M=\extM{r},
\]
naturally in $X$. Associativity of $\triangleleft_M$ is associativity of $\triangleleft$:
$(p\triangleleft_M q)\triangleleft_M s=p\triangleleft\ceil M\triangleleft q\triangleleft\ceil
M\triangleleft s=p\triangleleft_M(q\triangleleft_M s)$.
\end{proof}

\begin{proposition}[Non-unitality]\label{prop:nonunit}
The operation $\triangleleft_M$ has a two-sided unit up to ordinary container isomorphism if and only
if $M=\Id$.
\end{proposition}

\begin{proof}
If $M=\Id$ then $\ceil M=\y$ and $p\triangleleft_M q=p\triangleleft q$, with unit $\y$. Conversely,
let $e$ be a two-sided unit up to $\Cont$-isomorphism. Taking $q=\y$ in the left unit law,
$e\triangleleft_M\y=e\triangleleft\ceil M\cong\y$; applying the fully faithful $\ext{-}$ and
\eqref{eq:AAGmon} gives $\ext{e}\circ M\cong\Id$. Write $\ext{e}=\sum_{s\in S_e}\Set(E_s,-)$, so
$\sum_{s}\Set(E_s,M{-})\cong\Id$. Evaluate:
\begin{itemize}[leftmargin=1.4em]
\item At $X=1$: $\sum_s|M1|^{|E_s|}=1$. Each summand is $\ge1$ (as $|M1|\ge1$ by $\eta_1$), so there
  is a unique shape $s_0$ with $|M1|^{|E_{s_0}|}=1$; and $E_{s_0}=\emptyset$ is impossible
  ($\Set(\emptyset,M{-})=\mathrm{const}_1$ gives value $1\ne2$ at $X=2$). Hence $|M1|=1$ and
  $\Set(E_{s_0},M{-})\cong\Id$.
\item At $X=2$: $|M2|^{|E_{s_0}|}=2$. If $|E_{s_0}|\ge2$ this is $1$ or $\ge4$, never $2$; so
  $|E_{s_0}|=1$ and $|M2|=2$.
\end{itemize}
Thus $E_{s_0}\cong1$ and $\Set(1,M{-})=M\cong\Id$.
\end{proof}

\begin{remark}
The clean statement is a unit up to \emph{ordinary container} isomorphism, obtained by pushing the
unit law through the fully faithful $\ext{-}$. A unit up to the coarser $M\text{-}\Cont$
(Kleisli) isomorphism is a more delicate question we do not address. For applications the operative
fact is: $(\Cont,\triangleleft_M)$ is semigroupal, and monoidal only at $M=\Id$.
\end{remark}

\begin{remark}[Scope on morphisms]\label{rem:bifun}
Proposition~\ref{prop:suff} defines $\triangleleft_M$ on objects and makes $\extM{-}$ strong
semigroupal on the subcategory of \emph{pure} morphisms (those of the form $\eta\circ g$). Extending
$\triangleleft_M$ to a bifunctor on all of $M\text{-}\Cont$ so that $\extM{-}$ is strong semigroupal
requires transporting horizontal composites back along $\extM{-}$, which is free exactly when
$\extM{-}$ is full --- i.e.\ when $M$ is codense (Theorem~\ref{thm:2}). Thus the object-level closure
holds for every polynomial $M$, and upgrades to the full Kleisli category precisely when $M$ is
additionally codense. This Kleisli-bifunctoriality is the exact meeting point of
Theorems~\ref{thm:2} and \ref{thm:3}; we leave it open.
\end{remark}

\subsection{Necessity: the cardinality growth law}\label{sec:card}
Sufficiency looks like half of a biconditional, but the converse is subtler. We prove a
cardinality obstruction that already decides all the affine and commutative candidates, then isolate
the one remaining categorical gap.

\begin{proposition}[Cardinality growth law]\label{prop:cardlaw}
Suppose the essential image $\{\ext{r}\circ M:r\in\Cont\}$ is closed under $\circ$. Then for every
container $q$ there is a fixed nonnegative-integer polynomial $R_q$ with, writing $m(n)=|M(n)|$ and
$Q$ the cardinality polynomial of $\ext{q}$,
\[
  m\bigl(Q(m(n))\bigr)=R_q\bigl(m(n)\bigr)\qquad\text{for all finite }n.
\]
In particular, at $q=\y$ \textup(so $Q=\Id$\textup), $m(m(n))=R(m(n))$: the function $m$ agrees on
$\im(m)$ with a fixed nonnegative-integer polynomial.
\end{proposition}

\begin{proof}
The identity container $\y$ has $\extM{\y}=M$, and every $\extM{q}=\ext{q}\circ M$ lies in the image.
Closure applied to $p=\y$ and general $q$ gives, via Lemma~\ref{lem:A0},
$M\circ\ext{q}\circ M=\extM{\y}\circ\extM{q}\cong\ext{r_q}\circ M$ for some container $r_q$. Taking
cardinalities on a finite set of size $n$ yields $m(Q(m(n)))=R_q(m(n))$ with $R_q$ the (fixed)
cardinality polynomial of $r_q$. The case $q=\y$ is $Q=\Id$.
\end{proof}

\begin{corollary}[The affine/commutative folklore is false]\label{cor:folklore}
Nonempty powerset $\mathcal P^{+}$, powerset $\mathcal P$ and the distribution monad $\D$ all fail
composition, already at $p=q=\y$; while $\Maybe$ and the writer $E\times(-)$ compose. In particular
``composes $\iff$ commutative and affine'' is false in both directions.
\end{corollary}

\begin{proof}
For $\mathcal P^{+}$, $m(n)=2^{n}-1$ and $m(m(n))=2^{2^{n}-1}-1$ grows doubly exponentially in $n$,
hence super-polynomially in $v=m(n)$; no fixed polynomial $R$ satisfies $R(2^n-1)=2^{2^n-1}-1$ for
all $n$, so Proposition~\ref{prop:cardlaw} fails at $q=\y$. The same argument applies to $\mathcal P$
($m(n)=2^n$). For $\D$, the finite support map $\D\to\mathcal P^{+}$ is a monad morphism onto the
affine-commutative finite structure, reproducing the $\mathcal P^{+}$ growth on cardinalities, so $\D$
fails likewise. On the other side, $\Maybe=X+1$ and $E\times(-)$ are polynomial, so they compose by
Proposition~\ref{prop:suff}. Both $\mathcal P^{+}$ and $\D$ are affine and commutative yet fail;
$\Maybe$ and the writer are non-affine yet compose.
\end{proof}

\subsection{The categorical converse}\label{sec:converse}
Proposition~\ref{prop:cardlaw} is necessary but does not by itself pin $M$ as a functor: cardinalities
cannot separate $M$ from a polynomial functor with the same finite counts. The clean categorical
converse we want is

\begin{lemma}[the converse statement]\label{lem:N}
Consider the implication: if the essential image of $\extM{-}$ is closed under $\circ$ ---
equivalently, $M\circ N\circ M$ is of the form $\ext{r}\circ M$ for every polynomial $N$ --- then $M$
is polynomial.
\end{lemma}

We establish Lemma~\ref{lem:N} for four classes of $M$ (super-polynomial growth; analytic with
$M\emptyset=\emptyset$; bounded wreath depth; free commutative unital magmas), and we exhibit a monad
for which it \emph{fails} (the $C_3$ operad, \S\ref{sec:unital}). The frame is classical: over $\Set$,
\emph{$M$ is polynomial $\iff$ $M$ preserves connected limits} (wide pullbacks together with
cofiltered limits) $\iff$ $M$ is familially representable (Carboni--Johnstone; \cite[\S1]{GambinoKock},
\cite{NiuSpivak}). Two features cut the problem down.

\emph{First, $M$ is a monad, hence a retract of its square.}

\begin{lemma}[Retract reduction]\label{lem:R0}
For a monad $M$ on $\Set$: $M$ is polynomial $\iff$ $M\circ M$ preserves connected limits.
\end{lemma}

\begin{proof}
($\Rightarrow$) Polynomial functors are closed under composition \cite{GambinoKock}, so $M\circ M$ is
polynomial and preserves connected limits. ($\Leftarrow$) The right-unit law $\mu\circ\eta M=1_M$
exhibits $M$ as a \emph{split retract} of $M\circ M$ (section $\eta M$, retraction $\mu$).
Limit-preservation passes to retracts --- if the comparison map for $M\circ M$ is invertible then so
is the one for its retract $M$, by a diagram chase using naturality of comparison maps in the functor
variable. Hence $M$ preserves connected limits, so $M$ is polynomial.
\end{proof}

\emph{Second, the trailing $M$ cannot be cancelled, and this is now an exact statement rather than a
worry.} Closure at $p=q=\y$ yields only the single self-composite identity $M\circ M\cong\ext{r}\circ
M$, and that identity is categorically inert.

\begin{proposition}[The single instance is inert]\label{prop:inert}
If $M\circ M\cong P\circ M$ with $P$ polynomial and nonconstant, then for every connected-limit
diagram $D$, $M\circ M$ preserves $D$ $\iff$ $M$ preserves $D$. Hence the hypothesis $M\circ M\cong
\ext{r}\circ M$ carries no information about connected-limit preservation of $M$ beyond
Lemma~\ref{lem:R0}.
\end{proposition}

\begin{proof}
Naturally isomorphic functors preserve the same limits, so it suffices to treat $P\circ M$. The
comparison map for the composite factors as $\varphi_{PM}=\bigl(P(\lim MD)\cong\lim(P\circ MD)\bigr)
\circ P(\varphi_M)$, the left isomorphism holding because $P$, being polynomial, preserves the
connected limit $\lim(MD)$. A nonconstant polynomial $P=\sum_i(-)^{B_i}$ (some $B_i\ne\emptyset$)
reflects isomorphisms, so $\varphi_{PM}$ is invertible $\iff$ $\varphi_M$ is. Thus $M\circ M$ preserves
$D$ $\iff$ $M$ does.
\end{proof}

This is the precise form of the trailing-$M$ obstruction: precomposition $(-)\circ M$ is not
conservative, and the self-composite relates $M$ only to itself, so any proof of Lemma~\ref{lem:N}
must use either the \emph{full} family $\{M\circ N\circ M\cong\ext{r_N}\circ M:N\text{ polynomial}\}$
or a genuinely different invariant. We supply the latter --- a permutation-group invariant read off
the automorphisms of each structure --- and trace exactly how far it reaches, and where it stops.

\subsection{The stabilizer invariant and block-fixing rigidity}\label{sec:stab}
A natural isomorphism of $\Set$-endofunctors is $\mathrm{Aut}(X)$-equivariant at each $X$
(naturality at automorphisms $\sigma\colon X\to X$), so it preserves point-stabilizers.

\begin{lemma}[Stabilizer invariant]\label{lem:stab}
If $\Theta\colon F\cong G$ then $\mathrm{Stab}_{FX}(x)=\mathrm{Stab}_{GX}(\Theta_X x)\subseteq
\mathrm{Aut}(X)$ for all $X,x$. Hence the family of subgroups $\mathcal S(F):=\{\mathrm{Stab}_{FX}(x)\}$,
up to conjugacy, is an isomorphism invariant of $F$. Read on the \emph{each-label-once} components
below, it records one subgroup per structure and so is defined even when a component is infinite.
\end{lemma}

The master fact is that a \emph{polynomial} outer layer is rigid on positions. Call a structure
of an analytic functor \emph{each-label-once} on $[m]$ if its content is $[m]$ with every label used
exactly once (the multilinear part of the species component).

\begin{proposition}[Block-fixing rigidity of polynomial precomposition]\label{prop:meta}
Let $\ext r=\sum_{s}\Set(B_s,-)$ be polynomial and $M$ any endofunctor. On each each-label-once
component, every point-stabilizer of $(\ext r\circ M)[m]$ is a \emph{block-fixing product}
$\prod_{\beta\in\pi}\mathrm{Stab}_{M}(u_\beta)$ of single-$M$-element stabilizers, over the partition
$\pi$ of $[m]$ into the disjoint supports of the parts. In particular no two $M$-parts are ever
permuted.
\end{proposition}

\begin{proof}
An element of $(\ext r\circ M)[m]$ is $(s,\varphi\colon B_s\to M[\,\cdot\,])$ with total content
$[m]$. For $\sigma\in S_m$, $(\ext r\circ M)(\sigma)(s,\varphi)=(s,M\sigma\circ\varphi)$, which equals
$(s,\varphi)$ iff $M\sigma(\varphi(b))=\varphi(b)$ for \emph{every} $b\in B_s$: a function is fixed iff
fixed pointwise, there being no quotient on the positions $B_s$. So
$\mathrm{Stab}(s,\varphi)=\bigcap_{b}\mathrm{Stab}(\varphi(b))$. Each-label-once forces the parts
$\varphi(b)$ of nonempty content to have \emph{disjoint} supports $\beta$ partitioning $[m]$ (a shared
label would be used twice), and for disjoint blocks $\bigcap_\beta[\,\mathrm{Aut\ on\ }\beta\times
\mathrm{Sym}(\text{rest})\,]=\prod_\beta\mathrm{Aut\ on\ }\beta$.
\end{proof}

The each-label-once reading is essential and non-circular: on the \emph{raw} set $(\ext r\circ M)(X)$,
where labels may repeat, the stabilizer is an intersection over parts with possibly overlapping
supports --- a strictly larger class of subgroups. It is exactly the disjointness on the multilinear
component that collapses the intersection to a block-fixing product, and it is on that component that
Lemma~\ref{lem:stab} is compared. The first application recovers the universal analytic monad that
evades the cardinality law.

\begin{theorem}[The multiset monad does not compose]\label{thm:multiset}
Let $\mathbb M$ be the finite-multiset (free commutative monoid) monad, $\mathbb MX=\{$finite
multisets on $X\}$. Then $\mathbb M\circ\mathbb M\ncong P\circ\mathbb M$ for every polynomial $P$;
in particular $\mathbb M$ does not have composition-closed image.
\end{theorem}

\begin{proof}
A multiset $\nu\in\mathbb MX$ is a finitely-supported $\nu\colon X\to\mathbb N$, with
$\mathrm{Stab}_{\mathbb MX}(\nu)=\prod_k\mathrm{Sym}(\nu^{-1}(k))$ a \emph{Young subgroup} --- it
permutes points only \emph{within} the level-blocks of $\nu$. Intersections of Young subgroups are
Young, so by Proposition~\ref{prop:meta} every stabilizer of $(P\circ\mathbb M)[m]$ is a Young
subgroup. But for $X$ containing four distinct points $a,b,c,d$, the element
$m=\{\{a,b\},\{c,d\}\}\in\mathbb M\mathbb MX$ (content each-once) has
\[
  \mathrm{Stab}_{\mathbb M\mathbb MX}(m)=\bigl(\mathrm{Sym}\{a,b\}\times\mathrm{Sym}\{c,d\}\bigr)\rtimes
  \langle(a\,c)(b\,d)\rangle=S_2\wr S_2=D_4,\quad|D_4|=8,
\]
the within-block symmetries \emph{together with the block-swap} $(a\,c)(b\,d)$. $D_4$ contains
$(a\,c)(b\,d)$ but not the within-a-block transposition $(a\,c)$ (which would send $m$ to
$\{\{c,b\},\{a,d\}\}\ne m$), so $D_4$ is not a Young subgroup: a Young subgroup containing
$(a\,c)(b\,d)$ would place $a,c$ in one block and hence also contain $(a\,c)$. (Concretely the Young
subgroup orders in $S_4$ are $\{1,2,4,6,24\}$, and $8$ is not among them.) Thus $D_4\in\mathcal
S(\mathbb M\mathbb M)$ but $\mathcal S(P\mathbb M)$ consists of Young subgroups only, and
Lemma~\ref{lem:stab} gives $\mathbb M\mathbb M\ncong P\mathbb M$.
\end{proof}

The block-swap is an irreducibly second-order symmetry that a free (polynomial) outer layer cannot
manufacture: a single inner $\mathbb M$ supplies only within-block (Young) symmetry, whereas a
non-free \emph{outer} layer supplies between-equal-block (wreath) symmetry. In the species calculus
this is exactly the distinction between the labelled products in $B\bullet A$ (direct-product
stabilizers) and the wreath constituents of the self-substitution $A\bullet A$.

The multiset witness is a single wreath $S_2\wr S_2$ against Young subgroups; it reaches only
\emph{bounded} nesting. For monads whose structures nest to unbounded depth --- free magmas, free
operad-algebra monads, with automorphism towers $S_2\wr\cdots\wr S_2$ --- no fixed stabilizer suffices,
and the right invariant is the \emph{depth} of the wreath tower.

\subsection{The wreath-depth invariant}\label{sec:depth}
Recall the analytic/species dictionary (Joyal \cite{Joyal}; Bergeron--Labelle--Leroux \cite{BLL};
Gambino--Kock \cite{GambinoKock}). An \emph{analytic} functor is $\tilde A(X)=\sum_{n\ge0}A[n]\times_{S_n}
X^{n}$ for a \emph{species} $A$ (each $A[n]$ a finite $S_n$-set); the embedding $A\mapsto\tilde A$ is
fully faithful, so $\tilde A\cong\tilde B$ iff $A[n]\cong B[n]$ as $S_n$-sets for all $n$ (Joyal). Its
molecular decomposition is $A=\bigsqcup_l X^{d_l}/H_l$ with $H_l\le S_{d_l}$ the automorphism group of a
representative $A$-structure of arity $d_l$; $\tilde A$ is \emph{polynomial} iff $A$ is \emph{flat}
(every $H_l=1$, equivalently every $S_n$-action free), by familial representability
\cite[\S1]{GambinoKock}. Self-composition is species substitution, $\tilde A\circ\tilde A=\widetilde{A
\bullet A}$. By Proposition~\ref{prop:meta} every stabilizer of $(\ext r\circ\tilde A)[m]$ is a
block-fixing product of single-tree stabilizers, each a $\mathrm{Sym}(\beta)$-conjugate of some $H_l$;
we separate $\tilde A\circ\tilde A$ from $\ext r\circ\tilde A$ by a conjugacy invariant of subgroups
that block-fixing products cannot inflate.

\begin{definition}[Imprimitivity depth]\label{def:depth}
For a transitive $T\le\mathrm{Sym}(\Omega)$, a \emph{block system} is a $T$-invariant partition; let
$h_t(T)$ be the maximal length of a strictly refining chain of block systems from $\{\Omega\}$ to the
singletons (equivalently, of subgroups from a point-stabilizer up to $T$). For general
$K\le\mathrm{Sym}(\Omega)$, set $h(K)=\max\{h_t(K^O):O\text{ an orbit},\ |O|\ge2\}$, and $0$ if $K$ is
trivial, where $K^O$ is the transitive group $K$ induces on $O$. The \emph{wreath depth} of a species
is $\mathfrak h(A):=\sup_l h(H_l)$; $A$ has \emph{bounded wreath depth} iff $\mathfrak h(A)<\infty$.
\end{definition}

Both $h_t$ and $h$ are conjugacy invariants (conjugation carries block systems to block systems).
Sample values: $h_t(S_n)=1$ (primitive groups have no proper nontrivial block system),
$h_t(S_a\wr S_b)=2$, $h_t(S_2\wr S_2\wr S_2)=3$.

\begin{lemma}[Wreath superadditivity; product depth]\label{lem:wreath}
For transitive $T_1,T_2$ in the product action, $h_t(T_1\wr T_2)\ge h_t(T_1)+h_t(T_2)$. For a
block-fixing product $\prod_\beta G_\beta$ on pairwise disjoint supports,
$h(\prod_\beta G_\beta)=\max_\beta h(G_\beta)$.
\end{lemma}

\begin{proof}
For the wreath: the partition $\{\{j\}\times\Omega_1:j\in\Omega_2\}$ is invariant with quotient action
$T_2$; a longest chain of $T_2$ lifts to $h_t(T_2)$ strict steps above it, and a longest chain of
$T_1$ applied uniformly in every block gives $h_t(T_1)$ strict steps below it, concatenating to a chain
of length $\ge h_t(T_1)+h_t(T_2)$. For the product: the orbits of $\prod_\beta G_\beta$ are exactly the
orbits of the individual $G_\beta$ (each factor fixes the other supports), so
$h(\prod_\beta G_\beta)=\max_\beta h(G_\beta)$.
\end{proof}

\begin{theorem}[Wreath-depth converse]\label{thm:Sprime}
Let $M=\tilde A$ be an analytic monad of bounded wreath depth $W=\mathfrak h(A)<\infty$. If $A$ is
non-flat then $\tilde A\circ\tilde A\ncong\ext r\circ\tilde A$ for every polynomial $\ext r$; hence
closure of the essential image forces $M$ polynomial. This subsumes the bounded-degree case (degree
$\le d$ implies $\mathfrak h\le d-1<\infty$) and, at $W=1$, reproves Theorem~\ref{thm:multiset} and the
entire symmetric-power class ($\mathfrak h(\mathbb M)=1$ since each $H_n=S_n$ is primitive).
\end{theorem}

\begin{proof}
\emph{(Bound.)} By Proposition~\ref{prop:meta} and Lemma~\ref{lem:wreath}, every stabilizer of
$(\ext r\circ\tilde A)[m]$ is a block-fixing product of $\mathrm{Sym}(\beta)$-conjugates of constituent
groups $H_l$, so has depth $\le\max_\beta h(H_{l_\beta})\le\mathfrak h(A)=W$.
\emph{(Witness.)} Since $A$ is non-flat, some constituent $c$ acts nontrivially on an orbit
($h_t(H_c^{O_c})\ge1$); since $W<\infty$, the sup is attained by a constituent $b$ with an orbit
$h_t(H_b^{O_b})=W$. Fill every slot of $c$ with the same each-label-once $b$-structure on disjoint
labels: the resulting $w\in\widetilde{A\bullet A}$ has $\mathrm{Stab}(w)=H_b\wr H_c$, and restricting to
the invariant sub-orbit $O_c\times O_b$ gives the transitive sub-wreath $H_b^{O_b}\wr H_c^{O_c}$, of
depth $\ge h_t(H_c^{O_c})+h_t(H_b^{O_b})\ge1+W$ by Lemma~\ref{lem:wreath}. So $\tilde A\circ\tilde A$
realises depth $\ge W+1>W$, and Lemma~\ref{lem:stab} gives $\tilde A\circ\tilde A\ncong\ext r\circ\tilde A$.
\end{proof}

The invariant $h$ is blind precisely when $\mathfrak h(A)=\infty$: then $A$ already contains
arbitrarily deep wreath towers, $\tilde A\circ\tilde A$ produces no \emph{new} depth, and both sides
have infinite depth. That residual splits by $a_0=|M\emptyset|$: for $a_0=0$ an algebraic cancellation
reaches every depth (\S\ref{sec:cancel}); for $a_0>0$ the answer turns on symmetry
(\S\ref{sec:unital}).

\subsection{Cycle-index cancellation: the case $M\emptyset=\emptyset$}\label{sec:cancel}
For unbounded wreath depth the invariant $h$ cannot separate, but when $M\emptyset=\emptyset$ an
algebraic cancellation on cycle indices does, at any depth. The \emph{cycle-index series} of a species
$A$ is
\[
  Z_A=\sum_{n\ge0}\frac1{n!}\sum_{\sigma\in S_n}\mathrm{fix}(A[\sigma])\,p_\sigma\in
  R:=\mathbb Q[[p_1,p_2,\dots]],\qquad p_\sigma=\textstyle\prod_k p_k^{c_k(\sigma)}.
\]
The facts we need, beyond Joyal full faithfulness \textbf{(J1)} above. \textbf{(F$'$)} $\tilde A$ is
polynomial iff $Z_A\in\mathbb Q[[p_1]]$ (flatness in cycle-index form: a non-identity permutation fixing
a structure contributes a $p_\lambda$ with $\lambda\ne(1^n)$). \textbf{(C)} if $B[0]=\emptyset$ then
species substitution is plethysm on cycle indices, $Z_{A\bullet B}=Z_A\circ Z_B$, where $\circ$ is the
unique $\mathbb Q$-algebra endomorphism of $R$ with $p_k\circ H=H[p_i\mapsto p_{ki}]$. Writing
$a_0=|A[0]|=|M\emptyset|$ and $a_1=|A[1]|$, a monad has $a_1\ge1$ (the unit $\Id\Rightarrow\tilde A$
requires an $S_1$-fixed point of $A[1]$).

The engine is elementary.

\begin{lemma}[Plethysm right-cancellation]\label{lem:cancel}
Let $H\in R$ have zero constant term and nonzero linear coefficient, $H=a_1p_1+(\deg\ge2)$ with
$a_1\ne0$. Then $(-)\circ H\colon R\to R$ is injective: $F\circ H=G\circ H\Rightarrow F=G$.
\end{lemma}

\begin{proof}
$(-)\circ H$ is a $\mathbb Q$-algebra endomorphism, so it suffices that $D\circ H=0\Rightarrow D=0$.
Since $H$ has bottom degree $1$ with $H^{(1)}=a_1p_1$, plethysm gives $p_\lambda\circ
H=a_1^{\ell(\lambda)}p_\lambda+(\deg>|\lambda|)$. If $D\ne0$, let $m$ be its bottom degree; then
$(D\circ H)^{(m)}=\sum_{|\lambda|=m}c_\lambda\,a_1^{\ell(\lambda)}p_\lambda$, and linear independence of
$\{p_\lambda\}_{|\lambda|=m}$ together with $a_1\ne0$ forces every $c_\lambda=0$, contradicting
minimality.
\end{proof}

The hypothesis $a_1\ne0$ is sharp: with $a_1=0$ (e.g.\ $H=p_2+\tfrac12p_1^2$) plethysm collapses in low
degree. Both the injectivity and its failure at $a_1=0$ are computationally verified.

\begin{theorem}[Cancellation converse; $M\emptyset=\emptyset$]\label{thm:P}
Let $M=\tilde A$ be an analytic monad with $A[0]=\emptyset$. If the essential image of $\extM{-}$ is
closed under composition then $M$ is polynomial --- with no bound on degree or nesting depth.
\end{theorem}

\begin{proof}
Closure at $p=q=\y$ gives $M\circ M\cong\ext r\circ M$ for a polynomial $\ext r$. As $M\circ M$ is
analytic it is finitary, so $\ext r\circ\tilde A$ is finitary; and $\tilde AX\supseteq A[1]\times X$
(using $a_1\ge1$) is unbounded in $|X|$, so an infinite exponent $B_s$ in $\ext r=\sum_s(-)^{B_s}$
would make $(-)^{B_s}\circ\tilde A$ fail to preserve a filtered colimit --- forcing all $B_s$ finite.
Thus $\ext r=\tilde B$ for a \emph{flat} species $B$ (fact (F$'$)). Evaluating the isomorphism at $\emptyset$ and using $\tilde A\emptyset=A[0]=\emptyset$ gives
$B[0]=\tilde B\emptyset=\emptyset$. So $A\bullet A$ and $B\bullet A$ are finitary species with
$\widetilde{A\bullet A}=\tilde A\circ\tilde A$ and $\widetilde{B\bullet A}=\tilde B\circ\tilde A$; the
functor isomorphism and (J1) give $A\bullet A\cong B\bullet A$, hence equal cycle indices, and by (C)
\[
  Z_A\circ Z_A=Z_{A\bullet A}=Z_{B\bullet A}=Z_B\circ Z_A.
\]
Here $H:=Z_A$ has zero constant term ($a_0=0$) and $a_1\ge1$, so Lemma~\ref{lem:cancel} cancels $Z_A$
on the right: $Z_A=Z_B$. But $B$ is flat, so $Z_B\in\mathbb Q[[p_1]]$; hence $Z_A\in\mathbb Q[[p_1]]$,
i.e.\ $A$ is flat and $M$ is polynomial (F$'$).
\end{proof}

This is precisely the cancellation of the trailing $M$ that Proposition~\ref{prop:inert} forbade at the
level of limit preservation --- legitimate here because it operates on the \emph{faithful} cycle-index
invariant, which the single self-composite does determine. It reaches the monads no stabilizer count
can: the free commutative magma (structure counts $1,1,3,15,105,\dots=(2n-3)!!$, $a_0=0$, a balanced
$2^k$-leaf tree with automorphism group $S_2\wr\cdots\wr S_2$, $k$ factors) has unbounded nesting depth
yet falls to Theorem~\ref{thm:P} in one stroke.

\subsection{The unital corner: symmetry decides}\label{sec:unital}
After Theorems~\ref{thm:P} and \ref{thm:Sprime} the residual is: $M=\tilde A$ analytic, $a_0=|M\emptyset|$
nonempty-finite, and unbounded wreath depth $\mathfrak h(A)=\infty$. Cancellation needs $a_0=0$; the depth
invariant needs $\mathfrak h<\infty$; neither applies. This corner is inhabited, and --- the surprise of
this section --- it contains monads on \emph{both} sides of the converse.

Its canonical inhabitant is the \emph{free commutative unital magma} $U$: the free-algebra monad of one
commutative binary operation $\mu$ and a nullary unit $e$ absorbed by $\mu$ ($\mu(x,e)=x$). Its species
has $U[0]=\{e\}$ ($a_0=1$) and $U[n]=\{$commutative binary trees on $n$ leaves$\}$, $|U[n]|=(2n-3)!!$; the
balanced $2^k$-leaf tree has stabilizer $S_2\wr\cdots\wr S_2$ ($k$ factors), so $\mathfrak h(U)=\infty$
while every component is finite and $a_0=1$. It was the sharpest open case of the earlier form of this
theorem. We show necessity \emph{holds} at $U$ --- indeed for its whole symmetric family --- and then that
it \emph{fails} one operad over.

\subsubsection*{Full symmetry: necessity holds}
Call an operad $O$ \emph{fully symmetric unital} if every operation of arity $\ge2$ has full slot-symmetry
$H_g=S_{\mathrm{arity}}$, there is a nullary unit $e$ absorbed by every operation, the only unary operation
is the identity, and some operation has arity $\ge2$. Its free-algebra monad $M=\tilde A$ (with $A=O$) is
analytic, non-flat (some $H_g=S_{\ge2}\ne1$) hence non-polynomial, with $a_0\ge1$ and $\mathfrak h=\infty$.
Each such $O$ has a commutative binary unit-composite $\mu(x,y)=\omega(x,y,e,\dots,e)$ (delete $n-2$ slots
of an arity-$n$ operation by the unit; commutative by full symmetry). Write $\mathcal F(A)$ for the set of
support-indecomposable factors of automorphism groups of single $O$-trees (each-label-once). By
Proposition~\ref{prop:meta} and Lemma~\ref{lem:wreath}, every support-indecomposable factor of a
$(\ext r\circ M)[m]$ stabilizer lies in $\mathcal F(A)$; so it suffices to realise, in $(M\circ M)[m]$, a
factor \emph{outside} $\mathcal F(A)$.

\begin{lemma}[No rigid small tree; the escape factor]\label{lem:escape}
For a fully symmetric unital $O$: \textup{(R2)} every single $O$-tree on $\ge2$ labels has nontrivial
automorphism group; and consequently $\Delta_{C_2}:=\langle(0\,1)(2\,3)\rangle\notin\mathcal F(A)$.
\end{lemma}

\begin{proof}
(R2): in normal form no leaf is a unit (all absorbed), so a deepest internal node has only label-leaf
children and, by full symmetry, arity $\ge2$; a transposition of two of its leaves fixes the tree. For the
escape: under full symmetry $\mathrm{Aut}(t)=\prod_j(\mathrm{Aut}(s_j)\wr S_{m_j})$ over isomorphism
classes of children (the class-partition Young subgroup is a genuine direct product because $H_\rho=S_n$
contains every class-transposition independently), so every support-indecomposable factor is a wreath
$\Gamma=\mathrm{Aut}(s)\wr S_m$ ($m\ge2$) on $m\ell$ points, $\ell$ the label-count of $s$. If
$\Gamma\cong\Delta_{C_2}$ then $|\mathrm{Aut}(s)|^m\,m!=2$ and $m\ell=4$, forcing $m=2$,
$|\mathrm{Aut}(s)|=1$, $\ell=2$: a rigid $2$-label single tree, contradicting (R2).
\end{proof}

\begin{theorem}[Full symmetry gives necessity]\label{thm:family}
Let $M$ be the free-algebra monad of a fully symmetric unital operad. Then $M\circ M\ncong\ext r\circ M$
for every polynomial $\ext r$, so the essential image of $\extM{-}$ is not closed under composition; hence
$M$-containers compose $\iff$ $M$ is polynomial for every such $M$. This includes the free commutative
unital magma $U=U_2$ (the corner inhabitant) and its $n$-ary analogues $U_n$ for all $n\ge2$.
\end{theorem}

\begin{proof}
Let $V_i$ be the inner leaf on label $i$ and $E'$ the inner unit $e\in MX$ used as an \emph{outer} leaf
--- a rigid, non-absorbed peg (only the outer unit is absorbed at the outer level). Form
\[
  W_L=\mu(V_0,\mu(V_2,E')),\quad W_R=\mu(V_1,\mu(V_3,E')),\quad w=\mu(W_L,W_R)\in(M\circ M)[4].
\]
$W_L$ is rigid (its children $V_0$ and $\mu(V_2,E')$ are non-isomorphic, and inside the cherry $V_2\ne
E'$), likewise $W_R$, and $W_L\cong W_R$ via $0\mapsto1,\,2\mapsto3$; so
$\mathrm{Aut}(w)=\langle(0\,1)(2\,3)\rangle=\Delta_{C_2}$, whose only support-indecomposable factor is
itself. By Lemma~\ref{lem:escape} $\Delta_{C_2}\notin\mathcal F(A)$, so by Proposition~\ref{prop:meta} it
cannot be a stabilizer of $(\ext r\circ M)[4]$; Lemma~\ref{lem:stab} gives $M\circ M\ncong\ext r\circ M$.
\end{proof}

The mechanism is a \emph{correlated block-swap}: the unit builds a rigid block $W_L$ with no internal
symmetry, and the commutative outer $\mu$ swaps two equal copies, giving $(0\,1)(2\,3)$ --- a symmetry that
moves $0\leftrightarrow1$ \emph{only together with} $2\leftrightarrow3$. A block-fixing product cannot
correlate two blocks: it offers them independently ($S_2\times S_2$, order $4$) or fuses them into one
symmetric tree ($D_4$, order $8$); order $2$ is unreachable. The unit that inhabits the corner is exactly
what breaks composition there. (Verified for $U_3$: $\mathrm{Aut}(w)=\langle(0\,1)(2\,3)\rangle$, and an
exhaustive search over single $O$-trees on $4$ labels finds no $\Delta_{C_2}$ factor.)

\subsubsection*{Partial symmetry: necessity fails}
Full symmetry was load-bearing. Weaken it to \emph{cyclic} slot-symmetry and the converse collapses.

\begin{theorem}[The $C_3$ operad refutes necessity]\label{thm:C3}
Let $O$ be the operad freely generated by one ternary operation $\omega$ with cyclic slot-symmetry
$H_\omega=C_3=\langle(0\,1\,2)\rangle$ \textup{(}not $S_3$\textup{)} and an absorbed nullary unit $e$, and
let $M=\tilde A$ \textup{(}$A=O$\textup{)} be its free-algebra monad. Then $M$ is analytic and
non-polynomial \textup{(}non-flat, $H_\omega=C_3\ne1$\textup{)}, with $a_0=1$ and $\mathfrak h=\infty$ ---
it inhabits the same residual as $U$. Nevertheless the essential image of $\extM{-}$ \emph{is} closed under
composition: for all polynomial $p,q$,
\[
  \extM{p}\circ\extM{q}\;\cong\;\extM{p\triangleleft L},\qquad
  L=(\mathbb N,\ n\mapsto[n])\ \text{the list \textup{(}free-monoid\textup{)} container.}
\]
Hence composition-closure does \emph{not} imply polynomiality; the necessity direction of
Theorem~\ref{thm:3} is false without a symmetry hypothesis.
\end{theorem}

\begin{lemma}[Only three identical]\label{lem:c3}
For the $C_3$ operad, the only source of a nontrivial label-automorphism in a structure of $\ext p\circ
M\circ\ext q\circ M$ \textup{(}$p,q$ polynomial\textup{)} is a ternary $\omega$-node all three of whose
children are isomorphic \textup{(}contributing a $C_3$\textup{)}. The binary composite $\mu=\omega(x,y,e)$
is \emph{non-commutative} \textup{(}$C_3$ has no transposition\textup{)}, so two identical arguments are
never swapped, and the flat $p,q$ layers contribute only direct products. Consequently every
point-stabilizer is an iterated $C_3$-wreath assembled by direct products, and each of its
support-indecomposable factors is the automorphism group of a single $O$-tree.
\end{lemma}

\begin{proof}[Proof idea]
An automorphism of a rooted tree permutes each node's children within isomorphism classes using only that
node's slot-symmetry. Under $C_3$ a $3$-cycle fixes a configuration only when all three children coincide
(every $3$-cycle moves each slot to a non-identical one), and $\mu=\omega(-,-,e)$ has trivial residual
symmetry since $C_3$ has no transposition fixing the unit slot; flat layers have distinguishable positions.
So the only generator of label-symmetry is the three-identical $C_3$, the stabilizer is an iterated
$C_3$-wreath built by direct products, and each support-indecomposable factor is realised by a single
balanced $O$-tree.
\end{proof}

\begin{proof}[Proof of Theorem~\ref{thm:C3}]
Two analytic-type functors are isomorphic iff they carry the same multiset of \emph{molecule types}
$(n,[G])$ (arity and stabilizer conjugacy class), with multiplicity, on every component --- the
coefficientwise Joyal reconstruction, valid for infinite components. By Lemma~\ref{lem:c3} the molecule
types of $M\circ\ext q\circ M$ (any $q$ with a positive-arity shape) are exactly the Young products of
single-tree automorphism groups; the list composite $L\circ M$ realises the same set (a list of single
trees and units has a block-fixing-product stabilizer), and every such type recurs with multiplicity
$\aleph_0$ (pad by outer pegs, resp.\ appended units, without changing the stabilizer). Matching multisets
gives $M\circ\ext q\circ M\cong L\circ M$; whiskering by $\ext p$,
\[
  \ext p\circ M\circ\ext q\circ M\;\cong\;\ext p\circ L\circ M=(\ext p\circ\ext L)\circ M\cong
  \ext{p\triangleleft L}\circ M=\extM{p\triangleleft L}.
\]
The case $p=q=\y$ is $M\circ M\cong L\circ M$. As $M$ is non-flat it is non-polynomial, so the closed
image does not force polynomiality.
\end{proof}

The contrast with Theorem~\ref{thm:family} is exact --- the same construction seen from two operads. The
correlated-swap witness $w$ requires the outer $\mu$ to \emph{swap} its two equal children, which needs
$\mu$ commutative; under $C_3$, $\mu=\omega(-,-,e)$ is non-commutative, so $w$ has trivial automorphism
group and no escape occurs. (Verified: the molecule type-sets of $M\circ M$ and $L\circ M$ coincide for
$m\le4$; the same escape search that finds $\Delta_{C_2}$ for the \emph{fully symmetric} $S_3$-operad at
$m=4$ finds \emph{nothing} for $C_3$.)

\subsubsection*{The dividing line}
Both theorems turn on one bit --- the commutativity of the binary unit-composite
$\mu(x,y)=\omega(x,y,e,\dots)$.

\begin{center}\small
\renewcommand{\arraystretch}{1.25}
\begin{tabular}{lccc}
\hline
unital operad $O$ ($a_0>0$) & $\mu=\omega(x,y,e,\dots)$ & escape in $M\circ M$ & necessity\\
\hline
fully symmetric ($H_g=S$) & commutative & $\Delta_{C_2}\notin\mathcal F$ & \textbf{holds} (Thm~\ref{thm:family})\\
partially symmetric ($C_3$) & non-commutative & none & \textbf{fails} (Thm~\ref{thm:C3})\\
\hline
\end{tabular}
\end{center}

Commutative $\mu$ rigidifies a block with a unit peg and swaps two such blocks, manufacturing a correlated
symmetry no polynomial layer has; non-commutative $\mu$ manufactures nothing new, and every stabilizer is
already a single-tree automorphism group, so the image closes. This is the precise sense in which
\emph{polynomiality governs composition within the fully symmetric world}. What lies between the two rows
is open.

\begin{problem}[The non-symmetric residual]\label{prob:residual}
Characterise, among analytic monads with $M\emptyset$ nonempty-finite and unbounded wreath depth that are
\emph{not} fully symmetric, exactly which have composition-closed image. The $C_3$ operad populates the
``closed'' side and the fully symmetric unital magmas (Theorem~\ref{thm:family}) the ``not closed'' side;
a clean invariant separating them --- presumably graded by the symmetry of the binary unit-composites
available in $O$ --- would complete the necessity picture outside the symmetric world.
\end{problem}

Assembling sufficiency with the conditional converse:

\begin{theorem}[Composition and polynomiality; final form]\label{thm:3final}
\emph{Sufficiency \textup{(}unconditional\textup{)}:} if $M$ is polynomial then the essential image of
$\extM{-}$ is closed under composition, with $p\triangleleft_M q=p\triangleleft\ceil M\triangleleft q$
\textup{(}Proposition~\ref{prop:suff}\textup{)}. \emph{Necessity \textup{(}conditional\textup{)}:}
composition-closure forces $M$ polynomial for $M$ super-polynomial (Proposition~\ref{prop:cardlaw}), for
analytic $M$ with $M\emptyset=\emptyset$ (Theorem~\ref{thm:P}), for analytic $M$ of bounded wreath depth
(Theorem~\ref{thm:Sprime}), and for every fully symmetric unital magma monad (Theorem~\ref{thm:family},
including the corner case $U$); it \emph{fails} for the $C_3$ operad (Theorem~\ref{thm:C3}). The dividing
line is the commutativity of the binary unit-composite $\mu$. Consequently \emph{within the fully symmetric
world --- which contains every applied effect monad: multiset, distribution, powerset, symmetric powers ---
composition-closure $\iff$ flatness of the species $\iff$ $M$ polynomial}. The earlier reliance on an
unproved species-theoretic \emph{Plethysm Lemma} is removed (Theorem~\ref{thm:P}), and the formerly open
corner $U$ is resolved (necessity holds); the $C_3$ operad shows the biconditional cannot be extended past
the symmetry hypothesis.
\end{theorem}

\subsection{Independence}
\begin{proof}[Proof of Corollary~\ref{cor:tradeoff}]
The writer $E\times(-)$ is polynomial ($E\times(-)=\sum_{e\in E}(-)^1$), so it composes by
Proposition~\ref{prop:suff}; it is not codense ($\Phi(1)=|E|^{|E|}\ne|E|$ for $|E|\ge2$), so it is
not full by Theorem~\ref{thm:2}. The distribution monad $\D$ is affine and codense
(Example~\ref{ex:D}), so it is full by Theorem~\ref{thm:2}; it is not polynomial (it does not preserve
wide pullbacks --- joints are not determined by marginals), so it fails composition by
Corollary~\ref{cor:folklore}. Codensity and polynomiality are therefore logically independent
properties of $M$.
\end{proof}

\section{Conclusion}\label{sec:conc}

For a monad $M$ on $\Set$, the effectful container extension $\extM{-}=\ext{-}\circ M$ names
endofunctors faithfully always, names them \emph{fully} exactly when $M$ is codense, and has an image
closed under composition when $M$ is polynomial --- and, within the fully symmetric world that holds the
applied effect monads, \emph{only} when $M$ is polynomial. These are two different Kan-theoretic
invariants of the effect monad, and across the standard effects they pull in opposite directions:
probability is faithful but not composable, error is composable but not faithful. The folklore ``composes
iff commutative and affine'' is simply wrong --- $\D$ is affine and commutative but does not compose,
$\Maybe$ is neither yet does --- and the true obstruction to composition is a representability question
(flatness of a species), not a cohomology class or a distributive-law existence question. That
representability governs composition turns out to require a symmetry hypothesis: the $C_3$ operad, whose
non-polynomial monad nonetheless composes, marks the exact boundary.

For the compositional modelling of effectful agents this draws a sharp line. An agent that ``may
decline'' ($\Maybe$-positions) chains cleanly into other agents but cannot always be recovered from
its behaviour; an agent that ``answers probabilistically'' ($\D$-positions) is recoverable from its
behaviour but does not chain into a single probabilistic agent. One pays for declining in
faithfulness and for probability in composability, and the invariant that sets each price is now
explicit.

Three questions remain, each precisely located.
\begin{enumerate}[label=(\arabic*)]
\item \textbf{The composition converse, and its boundary (Problem~\ref{prob:residual}).} Does closure of
  the essential image under $\circ$ force $M$ to be polynomial? \emph{Conditionally yes, and in general
  no.} The converse holds for super-polynomial monads (the cardinality law), for analytic monads with
  $M\emptyset=\emptyset$ (plethysm right-cancellation, Lemma~\ref{lem:cancel}, at any wreath depth), for
  analytic monads of bounded wreath depth (the imprimitivity-depth invariant $h$,
  Definition~\ref{def:depth}), and for every fully symmetric unital magma monad (a correlated block-swap,
  Theorem~\ref{thm:family}) --- resolving the formerly open corner $U$ in favour of necessity. But the
  converse \emph{fails} for the $C_3$ operad (Theorem~\ref{thm:C3}): a non-polynomial monad whose image is
  nonetheless composition-closed, with $M\circ M\cong L\circ M$ for the list monad $L$. The trailing-$M$
  obstruction, sharpened to the ``no information'' statement of Proposition~\ref{prop:inert} at the level
  of limit preservation, is circumvented on the finer stabilizer and cycle-index invariants; and the exact
  dividing line is the commutativity of the binary unit-composite $\mu$. What remains open is the precise
  frontier between the two: which \emph{partially} symmetric unital analytic monads of unbounded wreath
  depth have composition-closed image (Problem~\ref{prob:residual})?
\item \textbf{Rigidity of $\D$ (Lemma~\ref{lem:Drigid}).} Full faithfulness of $\extM[\D]{-}$ is
  conditional on $\Nat(\D,\D)=\{\id\}$; the reduction of Example~\ref{ex:D} is unconditional. A clean
  proof of this rigidity fact would remove the one assumption in Theorem~\ref{thm:2}'s probabilistic
  case.
\item \textbf{Kleisli-bifunctoriality (Remark~\ref{rem:bifun}).} The composition product is defined on
  objects and is strong semigroupal on pure morphisms; it upgrades to the full Kleisli category
  exactly when $M$ is codense. Thus the two theorems meet: composition needs polynomiality to exist
  and codensity to become bifunctorial on effectful morphisms. The monads for which both hold --- the
  affine polynomial monads --- are worth isolating.
\end{enumerate}

\begin{thebibliography}{9}

\bibitem{AhmanUustaluDCC}
D.~Ahman and T.~Uustalu,
\emph{Directed containers as categories},
Electron.\ Proc.\ Theor.\ Comput.\ Sci.\ \textbf{207} (2016), 89--98.
arXiv:1604.01187.

\bibitem{BLL}
F.~Bergeron, G.~Labelle, and P.~Leroux,
\emph{Combinatorial Species and Tree-like Structures},
Encyclopedia of Mathematics and its Applications \textbf{67},
Cambridge University Press, 1998.

\bibitem{DJN}
A.~Dorta, C.~B.~Jarvis, and N.~Niu,
\emph{Monoidal structures on generalized polynomial categories},
arXiv:2305.05655.

\bibitem{GambinoKock}
N.~Gambino and J.~Kock,
\emph{Polynomial functors and polynomial monads},
Math.\ Proc.\ Cambridge Philos.\ Soc.\ \textbf{154} (2013), no.~1, 153--192.
arXiv:0906.4931.

\bibitem{Joyal}
A.~Joyal,
\emph{Foncteurs analytiques et esp\`eces de structures},
in \emph{Combinatoire \'enum\'erative} (Montr\'eal, 1985),
Lecture Notes in Math.\ \textbf{1234}, Springer, 1986, 126--159.

\bibitem{NiuSpivak}
N.~Niu and D.~I.~Spivak,
\emph{Polynomial Functors: A Mathematical Theory of Interaction},
book draft, 2023. arXiv:2312.00990.

\end{thebibliography}

\end{document}
